% Standalone supporting-information template inspired by the supplied article.
% Compile with: latexmk -pdf -interaction=nonstopmode main.tex
\documentclass[11pt,a4paper,oneside]{article}

\usepackage[T1]{fontenc}
\usepackage[scaled=0.96]{helvet}
\usepackage{newtxtext}
\renewcommand{\familydefault}{\sfdefault}
\usepackage[a4paper,left=22mm,right=22mm,top=23mm,bottom=23mm,headheight=22pt,headsep=9mm]{geometry}
\usepackage{microtype}
\usepackage{amsmath}
\usepackage[version=4]{mhchem}
\usepackage{chemfig}
\usepackage{xcolor}
\usepackage{booktabs}
\usepackage{tabularx}
\usepackage{siunitx}
\usepackage{graphicx}
\usepackage[font=small,labelfont=bf,justification=centering,singlelinecheck=false]{caption}
\usepackage{newfloat}
\DeclareFloatingEnvironment[fileext=los,name=Scheme,listname={LIST OF SUPPLEMENTARY SCHEMES}]{scheme}
\usepackage{titlesec}
\usepackage{tocloft}
\usepackage{fancyhdr}
\usepackage[hidelinks]{hyperref}

% Replace these fields with the matching manuscript metadata.
\newcommand{\journalname}{[JOURNAL NAME]}
\newcommand{\papertitle}{[Full title of the associated research article]}
\newcommand{\authors}{[Author One, Author Two, and Author Three]}
\newcommand{\correspondingauthor}{[Corresponding author and email address]}
\newcommand{\articleidentifier}{[Article DOI or manuscript identifier]}
\newcommand{\placeholder}[1]{\textit{[#1]}}

\setlength{\parindent}{0pt}
\setlength{\parskip}{5pt plus 1pt minus 1pt}
\setlength{\emergencystretch}{2em}
\linespread{1.12}
\setlength{\cftbeforesecskip}{2pt}
\setlength{\cftbeforesubsecskip}{0pt}
\renewcommand{\cftsecfont}{\sffamily\small}
\renewcommand{\cftsecpagefont}{\sffamily\small}
\renewcommand{\cftsubsecfont}{\sffamily\footnotesize}
\renewcommand{\cftsubsecpagefont}{\sffamily\footnotesize}
\renewcommand{\cftsecleader}{\cftdotfill{1.8}}
\renewcommand{\cftsubsecleader}{\cftdotfill{1.8}}
\renewcommand{\contentsname}{TABLE OF CONTENTS}
\renewcommand{\listfigurename}{LIST OF SUPPLEMENTARY FIGURES}
\renewcommand{\listtablename}{LIST OF SUPPLEMENTARY TABLES}

\titleformat{\section}
  {\large\sffamily\bfseries}
  {\thesection.}{0.55em}{}
\titleformat{\subsection}
  {\normalsize\sffamily\bfseries}
  {\thesubsection}{0.55em}{}
\titleformat{\subsubsection}
  {\normalsize\sffamily\bfseries}
  {\thesubsubsection}{0.55em}{}
\titlespacing*{\section}{0pt}{14pt plus 3pt minus 2pt}{6pt}
\titlespacing*{\subsection}{0pt}{10pt plus 2pt minus 1pt}{4pt}
\titlespacing*{\subsubsection}{0pt}{8pt plus 2pt minus 1pt}{3pt}

\pagestyle{fancy}
\fancyhf{}
\fancyhead[L]{\footnotesize\color{black!55}SUPPORTING INFORMATION}
\fancyhead[R]{\footnotesize\color{black!55}S\arabic{page}}
\renewcommand{\headrulewidth}{0.35pt}
\renewcommand{\headruleskip}{-5pt}
\renewcommand{\headrule}{\hbox to\headwidth{\color{black!25}\leaders\hrule height \headrulewidth\hfill}}
\renewcommand{\footrulewidth}{0pt}

\newcommand{\spectrumplaceholder}[4]{%
  \begin{minipage}[t]{\linewidth}
    \centering
    \IfFileExists{#3}{%
      \includegraphics[width=\linewidth,height=7.2cm,keepaspectratio]{#3}%
    }{%
      \begingroup
      \setlength{\fboxsep}{0pt}%
      \fbox{\parbox[c][7.2cm][c]{\dimexpr\linewidth-2\fboxrule\relax}{%
        \centering\color{black!55}%
        {\large\sffamily\bfseries #1 NMR}\par\vspace{8pt}%
        {\footnotesize #2}\par\vspace{12pt}%
        {\small Add the exported spectrum at\par\vspace{3pt}%
         \ttfamily\detokenize{#3}}%
      }}%
      \endgroup
    }%
    \captionof{figure}{#1 NMR spectrum of #2.}\label{fig:#4}%
  \end{minipage}%
}

\begin{document}

% The independent cover is deliberately excluded from the S-page sequence.
\pagenumbering{alph}
\thispagestyle{empty}
\begin{titlepage}
  \centering
  \vspace*{18mm}
  {\color{black!25}\rule{0.82\textwidth}{0.6pt}}\par
  \vspace{7mm}
  {\large\sffamily\bfseries SUPPORTING INFORMATION\par}
  \vspace{3mm}
  {\normalsize\color{black!60}\journalname\par}

  \vspace{30mm}
  {\color{black!25}\rule{0.4\textwidth}{0.45pt}\par}
  \vspace{13mm}
  {\fontsize{22}{28}\selectfont\bfseries\rmfamily\papertitle\par}
  \vspace{9mm}
  {\large\rmfamily\itshape\authors\par}
  \vspace{13mm}
  {\color{black!25}\rule{0.4\textwidth}{0.45pt}\par}

  \vfill
  {\small\color{black!65}\correspondingauthor\par}
  \vspace{3mm}
  {\footnotesize\color{black!55}\articleidentifier\par}
  \vspace{18mm}
\end{titlepage}

\pagenumbering{arabic}
\setcounter{page}{1}
\setcounter{tocdepth}{2}

\tableofcontents
\vspace{8pt}
\listoftables
\vspace{8pt}
\listofschemes
\vspace{8pt}
\listoffigures
\clearpage

\section{General Information}

\subsection{Materials and Methods}
\placeholder{Describe the source, grade, and handling of reagents and solvents.}

\placeholder{Describe the atmosphere, glassware preparation, and reaction setup where relevant.}

\subsection{Chromatography and Instrumentation}
\placeholder{List chromatography media, analytical instruments, field strengths, solvents, and measurement conditions. State how reported values and units were obtained.}

\subsection{Data and Compound Nomenclature}
\placeholder{Describe the software and conventions used to assign compound names, spectral data, and analytical results. Add repository or accession links for deposited data, where applicable.}

\section{Experimental Data}

\subsection{Optimization of Reaction Conditions}
Table~\ref{tab:optimization} provides a reusable format for reaction optimization. Replace every bracketed field with verified measurements and add explanatory footnotes as needed.

\begin{table}[htbp]
  \centering
  \caption{Optimization of \placeholder{reaction or experimental parameter}.}
  \label{tab:optimization}
  \small
  \begin{tabularx}{\linewidth}{@{}c X c c@{}}
    \toprule
    Entry & \placeholder{Condition or reagent} & \placeholder{Solvent} & Yield [\%] \\
    \midrule
    1 & \placeholder{Condition A} & \placeholder{Solvent A} & \placeholder{value} \\
    2 & \placeholder{Condition B} & \placeholder{Solvent B} & \placeholder{value} \\
    3 & \placeholder{Condition C} & \placeholder{Solvent C} & \placeholder{value} \\
    \bottomrule
  \end{tabularx}

  \vspace{4pt}
  \begin{minipage}{0.96\linewidth}
    \footnotesize
    \textit{Notes:} \placeholder{State equivalents, scale, assay method, replicate count, and the source of each yield.}
  \end{minipage}
\end{table}

\subsection{General Procedure}
\subsubsection{General Procedure A}
\label{sec:procedure-a}
\placeholder{Describe a representative procedure in the past tense. Specify substrate and reagent amounts, equivalents, solvent, concentration, temperature, time, work-up, purification, product mass, and isolated yield.}

\subsubsection{General Procedure B}
\label{sec:procedure-b}
\placeholder{Add another procedure only when it materially differs from the preceding procedure.}

\section{Compound Characterization}

\subsection{Product A (1)}
\label{sec:product-a}
\placeholder{Insert the full compound name and identifier used consistently in the manuscript, schemes, spectra, and raw data.}

\begin{scheme}[htbp]
  \centering
  \setchemfig{atom sep=1.8em}
  \chemfig{R-OH}
  \qquad\ensuremath{\longrightarrow}\qquad
  \chemfig{R-O-R'}
  \caption{Illustrative reaction layout. Replace the generic structures with the corresponding publication-ready chemical structures.}
  \label{sch:representative}
\end{scheme}

\paragraph{Preparation.}
\placeholder{State the procedure reference, starting material, scale, purification, appearance, isolated mass, and yield.}

\paragraph{Analytical data.}
\textbf{TLC:} \placeholder{$R_f$, solvent system, and visualization method.}

\textbf{NMR spectroscopy:} \placeholder{Solvent and instrument.}

\textsuperscript{1}H NMR (\placeholder{frequency} MHz, \ce{CDCl3}):
\placeholder{chemical shift, multiplicity, coupling constants, integrations, and assignments}; ppm.

\textsuperscript{13}C NMR (\placeholder{frequency} MHz, \ce{CDCl3}):
\placeholder{chemical shifts and assignments}; ppm.

\textbf{HRMS} (\placeholder{ionization mode}):
\placeholder{calculated $m/z$ for the assigned ion}; found:
\placeholder{measured $m/z$}.

\textbf{IR} (\placeholder{sampling method}):
\placeholder{diagnostic absorption maxima}; cm$^{-1}$.

\clearpage
\section{NMR Spectra}
Replace the example filenames below with exported, publication-ready spectrum files. The framed placeholders are included only to preserve the intended layout when an image has not yet been added.

\subsection*{Product A (1)}
\begin{figure}[htbp]
  \centering
  \spectrumplaceholder{\textsuperscript{1}H}{Product A (1)}{figures/1H-1.pdf}{nmr-1h-product-a}
  \par\vspace{8mm}
  \spectrumplaceholder{\textsuperscript{13}C}{Product A (1)}{figures/13C-1.pdf}{nmr-13c-product-a}
\end{figure}

\clearpage
\subsection*{Product B (2)}
\begin{figure}[htbp]
  \centering
  \spectrumplaceholder{\textsuperscript{1}H}{Product B (2)}{figures/1H-2.pdf}{nmr-1h-product-b}
  \par\vspace{8mm}
  \spectrumplaceholder{\textsuperscript{13}C}{Product B (2)}{figures/13C-2.pdf}{nmr-13c-product-b}
\end{figure}

\section{References}
\begingroup
\renewcommand{\labelenumi}{[\arabic{enumi}]}
\begin{enumerate}
  \item
  \placeholder{Author(s). Article title. \textit{Journal} Year, Volume, page range. DOI or persistent identifier.}
\end{enumerate}
\endgroup

\end{document}
